\documentclass[11pt]{article}
% Preámbulo compartido para los modelos de parcial ACN
\usepackage[utf8]{inputenc}
\usepackage[T1]{fontenc}
\usepackage[spanish]{babel}
\usepackage{amsmath,amssymb}
\usepackage[margin=2.3cm]{geometry}
\usepackage{enumitem}
\usepackage{fancyhdr}
\usepackage{array}

\pagestyle{fancy}
\fancyhf{}
\lhead{Aplicaciones Computacionales en Negocios --- UTDT}
\rhead{\thepage}
\renewcommand{\headrulewidth}{0.4pt}

% Encabezado tipo examen
\newcommand{\encabezado}[1]{%
  \begin{center}
    {\Large\bfseries Examen parcial}\\[2pt]
    {\large Aplicaciones Computacionales en Negocios --- UTDT}\\[2pt]
    {\normalsize #1}\\[8pt]
  \end{center}
  \noindent\fbox{\parbox{\linewidth}{\centering\bfseries RESPONDER EXCLUSIVAMENTE DEBAJO DEL ESPACIO PROVISTO BAJO CADA PREGUNTA}}
  \vspace{6pt}

  \noindent Nombre y apellido: \underline{\hspace{7cm}} \hfill Legajo: \underline{\hspace{3.5cm}}
  \vspace{6pt}
  \hrule
  \vspace{10pt}
}

% Puntaje en negrita
\newcommand{\pts}[1]{\textbf{[#1 puntos]}\ }

% Espacio de respuesta
\newcommand{\espacio}[1]{\vspace{#1}}


\begin{document}
\encabezado{Modelo 1}

% ------------------------------------------------------------------
\noindent\textbf{1)} A partir de la simulación de un modelo estocástico obtengo $N$ observaciones independientes de una variable $Y$. Con los primeros $2000$ caminos estimo una media de $150$ y una desviación estándar de $60$.

\medskip
\pts{20} ¿Cuántas observaciones (caminos) debo simular en total para que el error de estimación de la media de $Y$ sea menor a $0.4$? Explicá qué resultado teórico usás y cómo escala el error con el número de caminos.

\espacio{4.5cm}

\pts{10} Si en lugar de estimar la media quiero estimar la probabilidad de un evento raro con $\hat p \ll 1$, ¿por qué se vuelve más difícil (necesito más caminos) alcanzar un error relativo chico? Justificá con la fórmula del error de un estimador de probabilidad.

\espacio{4cm}

% ------------------------------------------------------------------
\noindent\textbf{2)} Considero dos procesos que evolucionan simultáneamente en el tiempo:
\[
X_{i+1} = X_i + \mu\,\Delta, \qquad X_0 = 5
\]
\[
Y_{i+1} = Y_i - \kappa\,(Y_i - X_i)\,\Delta + \sigma\sqrt{\Delta}\,Z_i, \qquad Y_0 = 0
\]
donde los $Z_i$ son shocks normales i.i.d.\ con media cero y desviación estándar $1$. Los parámetros son $\Delta = 1$, $\mu = 0.5$, $\sigma = 1$.

\medskip
\pts{10} Hacé un gráfico esquemático con $3$ caminos típicos de $X$ e $Y$ entre $i=0$ e $i=10$ si $\kappa = 0$. Describí en palabras el comportamiento de $Y$ respecto de $X$.

\espacio{5cm}

\pts{15} Hacé un gráfico esquemático con $3$ caminos típicos de $X$ e $Y$ entre $i=0$ e $i=10$ si $\kappa = 1$. ¿Qué rol cumple el término $-\kappa(Y_i - X_i)\Delta$? ¿Qué pasa cuando $\kappa$ es grande?

\espacio{5cm}

% ------------------------------------------------------------------
\noindent\textbf{3)} Quiero generar muestras de clientes de un banco, caracterizados por su segmento de riesgo y por el monto de un préstamo solicitado. La probabilidad de que un cliente pertenezca al segmento $j$ es $p_j$, para $j = 0,1,2,3,4$ (las probabilidades suman $1$). El monto solicitado por un cliente del segmento $j$ es aleatorio y uniforme entre $j\cdot 20000$ y $(j+1)\cdot 20000$ dólares. Suponé que tenés disponible una función para generar variables uniformes $[0,1]$.

\medskip
\pts{25} Proponé un algoritmo de simulación para generar el segmento y el monto solicitado de un cliente. Describí de manera precisa (pseudocódigo y palabras) cómo generás cada una de las variables aleatorias necesarias.

\espacio{6cm}

\pts{10} Describí con precisión cómo estimar, vía la simulación de un $N$ grande de clientes, el valor esperado del monto solicitado por los clientes de los segmentos $0$ y $1$ únicamente.

\espacio{4.5cm}

\pts{10} Describí con precisión cómo calcular el error de tu estimador a partir de la muestra simulada, y qué nuevo tamaño de muestra necesitarías si quisieras reducir ese error a la mitad.

\espacio{4cm}

\end{document}
